% =================================================================
% CONFIGURAÇÕES DO TEMPLATE (ESQUELETO)
% Centralize aqui customizações do documento.
% Como usar:
%  - Preencha os placeholders abaixo (título, autor, etc.)
%  - Adicione/remova pacotes conforme a necessidade do seu trabalho
%  - Mantenha overrides aqui, sem editar config/sbc-template.sty
% =================================================================

% Pacotes essenciais
\usepackage{config/sbc-template}
\usepackage[T1]{fontenc}   % melhor hifenização/acentuação com pdflatex
\usepackage[utf8]{inputenc}
\usepackage{microtype}     % microtipografia (kerning/hifenização)
% Idiomas: en + pt-BR (compatibilidade com babel antigo/novo)
\IfFileExists{english.ldf}{%
  % Babel legado (.ldf)
  \IfFileExists{brazilian.ldf}{%
    \usepackage[english,brazilian]{babel}%
  }{%
    \IfFileExists{portuguese.ldf}{%
      \usepackage[english,portuguese]{babel}%
    }{%
      \usepackage[english]{babel}%
    }%
  }%
}{%
  % Babel novo (babelprovide)
  \usepackage{babel}%
  \IfFileExists{babel-brazilian.tex}{%
    \babelprovide[main]{brazilian}%
    \babelprovide{english}%
  }{%
    \IfFileExists{babel-portuguese.tex}{%
      \babelprovide[main]{portuguese}%
      \babelprovide{english}%
    }{%
      \babelprovide[main]{english}%
    }%
  }%
}
\usepackage{csquotes}      % citações tipográficas
\usepackage{graphicx,url}
\usepackage{hyperref}
\graphicspath{{assets/}}   % pasta padrão para imagens

% Pacotes adicionais (habilite conforme necessário)
\usepackage{listings}           % listagens de código (opcional)
\usepackage{xcolor}             % cores (necessário para listings)
\usepackage{amsmath,amssymb}    % matemática
\usepackage{booktabs}           % tabelas profissionais (\toprule, etc.)
\usepackage{tabularx}           % tabelas largas com colunas X
\usepackage{array}              % colunas avançadas em tabelas
\usepackage{enumerate}          % listas enumeradas flexíveis
% Glossário e índice (opcionais):
% \usepackage[acronym]{glossaries} % requer makeglossaries; ver README
% \makeglossaries
% \usepackage{imakeidx}           % índice remissivo
% \makeindex
% BibLaTeX (opção B2):
% \usepackage[backend=biber,style=authoryear,doi=true,isbn=false,url=true]{biblatex}
% \addbibresource{bibliography/sbc-template.bib}
% Minted (opcional, para código com realce de sintaxe via Pygments):
% Requer: \usepackage{minted} e compilar com TEXFLAGS="-shell-escape".
% \usepackage{minted}

% Estilo de código (exemplo; ajuste ou remova)
\definecolor{codegray}{RGB}{248,248,248}
\lstdefinestyle{code}{
    backgroundcolor=\color{codegray},
    basicstyle=\footnotesize\ttfamily,
    numbers=left,
    numbersep=5pt,
    frame=single,
}

% Hyperlinks (preenche metadados do PDF)
\hypersetup{
    colorlinks=true,
    linkcolor=black,
    urlcolor=blue,
    citecolor=blue,
    pdftitle={[Título do Trabalho]},
    pdfauthor={[Nome do Autor]},
}

% Comandos de placeholders (consumidos em \title, \author e \address)
\newcommand{\DocumentTitle}{[Título do Trabalho]}
\newcommand{\AuthorName}{[Nome do Autor]}
\newcommand{\AuthorInstitution}{[Departamento / Curso] -- [Instituição]}
\newcommand{\AuthorLocation}{[Cidade] -- [UF] -- [País]}
\newcommand{\AuthorEmail}{[email@exemplo.com]}

% Metadados padrão (para projetos que usam este arquivo como entrada principal)
\title{\DocumentTitle}
\author{\AuthorName\inst{1}}
\address{\AuthorInstitution\\
  \AuthorLocation\\
  \email{\AuthorEmail}
}

% Configurações gerais
\sloppy
\setlength{\parskip}{6pt}
